\documentclass{article}

\usepackage{xcolor}
\usepackage[utf8]{inputenc}
\usepackage{amsmath}
\usepackage{enumerate}
\usepackage{microtype}
\usepackage[hyphens]{url}
\usepackage{hyperref}
\usepackage[top=1.5cm]{geometry} % Page margins

\title{Rúbrica Proyecto Final - Compresión de archivos de texto}
\author{Estructuras de Datos \\ Facultad de Ingenierías - Universidad Tecnológica de Pereira \\ Ingeniería de Sistemas y Computación}
\date{}

\begin{document}
\maketitle
\thispagestyle{empty}

\section{Requerimientos}
\begin{itemize}
    \item Implementar el algoritmo de \textbf{Huffman Coding} para la compresión de un archivo con extensión .txt a un archivo con extensión .bin y extracción (descompresión) de datos textuales a partir de un archivo con formato .bin.
    \item Implementar el algoritmo de compresión \textbf{LZ77} para la compresión de un archivo con extensión .txt a un archivo con extensión .bin y extracción (descompresión) correspondiente.
    \item Implementar un sistema integrado que combine ambos algoritmos de manera secuencial . Es decir, dado un archivo de texto:
    \begin{itemize}
        \item Aplicar inicialmente el algoritmo LZ77.
        \item Posteriormente, aplicar Huffman Coding sobre el resultado.
        \item Permitir la extracción de un archivo comprimido que haya sido producido a partir de este método
    \end{itemize}
    \item Para cada uno de los numerales anteriores se debe garantizar la reconstrucción \textbf{sin pérdidas (lossless)} de la información original al extraer.
    \item Implementar funcionalidades para evaluar el desempeño de cada uno de los anteriores numerales, considerando los siguientes criterios:
    \begin{itemize}
        \item Tasa de compresión.
        \item Tiempo de ejecución.
    \end{itemize}
    \item Elaboración de un reporte de resultados que incluya los resultados recolectados en el numeral anterior así como sus conclusiones con relación a estos y al proceso de desarrollo.
    \item El proyecto debe incluir una \textbf{interfaz basada en terminal} desde la cual permita:
    \begin{itemize}
        \item Elegir la tarea de compresión/extracción a ejecutar (Compresión/Extracción).
        \item Seleccionar el archivo sobre el cual se va a realizar el proceso.
    \end{itemize}
    \item El proyecto debe ser presentado y sustentado de manera grupal durante la última semana de clase.
    \item La implementación debe realizarse en lenguaje C/C++.
\end{itemize}


\textbf{Nota: NO} se requiere que se maneje persistencia entre ejecuciones del programa; los procesos de compresión y extracción de archivo se van a realizar en una sola instancia de ejecución.

\section{Bonus} 
\begin{itemize}
    \item Implementación del algoritmo \textbf{LZMA} (Lempel-Ziv-Markov chain algorithm), incorporando el uso de \textbf{cadenas de Markov}, como en el caso de uso de 7Z. Se valorará la integración con los requerimientos anteriores. \textbf{(+2.5)}.
    \item La implementación de herramientas adicionales que demuestren creatividad y entendimiento del problema y alcance del proyecto, así como mejoras al mismo serán consideradas como bonificación proporcionalmente.
\end{itemize}
\section{Lecturas recomendadas} 
\subsection{Huffman Coding}
\begin{itemize}
    \item Moffat, A. (2019). Huffman coding. ACM Computing Surveys (CSUR), 52(4), 1-35.
    \item \url{https://www.w3schools.com/dsa/dsa_ref_huffman_coding.php}
\end{itemize}
\subsection{LZ77}
\begin{itemize}
    \item \url{https://learn.microsoft.com/en-us/openspecs/windows_protocols/ms-wusp/fb98aa28-5cd7-407f-8869-a6cef1ff1ccb}
\end{itemize}
\subsection{LZMA}
\begin{itemize}
    \item \url{https://www.7-zip.org/sdk.html}
\end{itemize}

\end{document}